% META: {"id":"1SPE-DERGLOBAL-COURS-C1","chapitre":"1SPE-DERIVATION-GLOBAL","type_objet":"cours","capacites_codes":["C1"],"status":"generated"}

\section{Dérivées des fonctions de référence}

\margeAppui{Ce tableau est à connaître par cœur : il constitue la « boîte à outils » de toute la dérivation.}

\definition[Fonction dérivée]{Soit $f$ une fonction dérivable sur un intervalle $I$. La \textbf{fonction dérivée} de $f$, notée $f'$, est la fonction qui, à tout $x$ de $I$, associe le nombre dérivé $f'(x)$.}

\propriete[Tableau des dérivées de référence]{
La deuxième colonne indique où la fonction est dérivable. Ce domaine peut être
plus petit que son domaine de définition : la fonction racine carrée est
définie sur $[0\,;\,+\infty[$, mais elle n'est dérivable que sur $]0\,;\,+\infty[$.
\[
\begin{array}{|c|c|c|}
\hline
\textbf{Fonction } f & \textbf{Dérivabilité} & \textbf{Dérivée } f' \\
\hline
f(x) = k \quad (k \in \mathbb{R}) & \mathbb{R} & f'(x) = 0 \\
\hline
f(x) = x & \mathbb{R} & f'(x) = 1 \\
\hline
f(x) = x^n \quad (n \in \mathbb{N},\, n \geqslant 2) & \mathbb{R} & f'(x) = nx^{n-1} \\
\hline
f(x) = x^n \quad (n \in \mathbb{Z},\, n \leqslant -1) & \mathbb{R} \setminus \{0\} & f'(x) = nx^{n-1} \\
\hline
f(x) = \dfrac{1}{x} & \mathbb{R} \setminus \{0\} & f'(x) = -\dfrac{1}{x^2} \\[6pt]
\hline
f(x) = \sqrt{x} & ]0\,;\,+\infty[ & f'(x) = \dfrac{1}{2\sqrt{x}} \\[6pt]
\hline
\end{array}
\]
}

\exemple{
Pour $f(x) = x^5$, on identifie $n = 5$. La formule $f'(x) = nx^{n-1}$ donne :
\[
f'(x) = 5x^4.
\]
}

\exemple{
Pour $g(x) = x^3$, on a $n = 3$, donc :
\[
g'(x) = 3x^2.
\]
Par exemple, $g'(2) = 3 \times 4 = 12$. La tangente à la courbe de $g$ au point d'abscisse $2$ a pour pente $12$.
}

\exemple{
Pour $h(x) = \dfrac{1}{x}$, définie sur $\mathbb{R} \setminus \{0\}$ :
\[
h'(x) = -\frac{1}{x^2}.
\]
La dérivée est toujours \emph{strictement négative}, ce qui s'interprétera comme une décroissance de $h$ sur chacun de ses intervalles de définition.
}

\exemple{
Pour $\varphi(x) = \sqrt{x}$, définie sur $[0\,;\,+\infty[$ et dérivable sur $]0\,;\,+\infty[$ :
\[
\varphi'(x) = \frac{1}{2\sqrt{x}}.
\]
En $x = 4$ : $\varphi'(4) = \dfrac{1}{2\sqrt{4}} = \dfrac{1}{4}$.
}

% [ROOT-ZERO-PROOF]
\contreexemple{
La fonction $\varphi(x)=\sqrt{x}$ n'est pas dérivable en $0$. Pour $h>0$,
son taux de variation entre $0$ et $h$ vaut
\[
\frac{\sqrt{h}-\sqrt{0}}{h-0}=\frac{1}{\sqrt{h}}.
\]
Quand $h$ se rapproche de $0$ par valeurs positives, ce taux dépasse toute
borne fixée : il ne tend vers aucun réel. Il n'existe donc pas de nombre dérivé
en $0$. Graphiquement, la courbe possède une demi-tangente verticale en $(0\,;\,0)$ ;
on ne lui attribue pas une pente réelle.
}
% [/ROOT-ZERO-PROOF]

\contreexemple{
La fonction valeur absolue $f(x) = |x|$ n'est pas dérivable en $0$.

En effet, le taux de variation à gauche et à droite de $0$ donne :
\[
\lim_{h \to 0^-} \frac{|0+h| - |0|}{h} = \lim_{h \to 0^-} \frac{-h}{h} = -1 \qquad \text{(dérivée à gauche)}
\]
\[
\lim_{h \to 0^+} \frac{|0+h| - |0|}{h} = \lim_{h \to 0^+} \frac{h}{h} = 1 \qquad \text{(dérivée à droite)}
\]
Les limites à gauche et à droite sont différentes ($-1
\neq 1$) : le nombre dérivé n'existe pas en $0$.

\emph{Interprétation graphique :} la courbe de $|x|$ présente un « angle pointu » (un point anguleux) en $x = 0$. Il n'y a pas de tangente unique en ce point.
}

\margeAppui{Voici deux exemples de non-dérivabilité : une demi-tangente verticale pour $\sqrt{x}$ en $0$, et un point anguleux pour $|x|$ en $0$. Dans ces deux exemples, la fonction est continue en ce point.}

% BEGIN-VERIFY
% from sympy import symbols, Abs, sqrt, simplify, diff, limit, oo
% x = symbols('x', real=True)
% h = symbols('h', positive=True)
% assert simplify((sqrt(h)-sqrt(0))/h - 1/sqrt(h)) == 0
% assert limit(sqrt(h)/h, h, 0, dir='+') == oo
% # Identités générales des taux à droite et à gauche, avec h positif.
% assert simplify(Abs(h)/h) == 1
% assert simplify(Abs(-h)/(-h)) == -1
% assert diff(x, x) == 1
% assert diff(-x, x) == -1
% END-VERIFY

\erreurFrequente{
Écrire $(x^n)' = nx^n$ en oubliant de diminuer l'exposant d'une unité. La bonne formule est $(x^n)' = nx^{n-1}$ : on multiplie par l'exposant et on soustrait $1$.
}

\erreurFrequente{
Écrire $\left(\dfrac{1}{x}\right)' = \dfrac{1}{x^2}$ en oubliant le signe moins. La dérivée de $\dfrac{1}{x}$ est $-\dfrac{1}{x^2}$.
}

\erreurFrequente{
La racine carrée $\sqrt{x}$ n'est pas $\dfrac{x}{2}$. Pour $x>0$, sa dérivée
est $\dfrac{1}{2\sqrt{x}}$, et non $\dfrac{1}{2}$.
}

\margeAppui{Le cas $n=1$ est la fonction identité, de dérivée $1$ en tout réel.
Le cas $n=0$ donne $x^0=1$ pour $x\neq0$ : sa dérivée est $0$ sur ce domaine.
La fonction constante $1$ est, elle, définie et dérivable sur tout $\mathbb R$,
sans utiliser $0^0$ ni multiplier $0$ par une expression non définie.}

\subsection*{Deux démonstrations à connaître}

Pour la fonction carré, pour tout réel $a$ et tout $h\neq0$,
\[
\frac{(a+h)^2-a^2}{h}=2a+h.
\]
Ce taux tend vers $2a$ lorsque $h$ tend vers $0$ : la dérivée de la fonction
carré en $a$ est donc $2a$.

Pour la fonction inverse, fixons $a\neq0$. Pour $h\neq0$ assez proche de $0$
pour que $a+h\neq0$,
\[
\frac{\frac{1}{a+h}-\frac{1}{a}}{h}
=-\frac{1}{a(a+h)}.
\]
Lorsque $h$ tend vers $0$, ce taux tend vers $-\frac{1}{a^2}$.
Cela démontre la formule sur chacun des deux intervalles où la fonction
inverse est définie, pour les valeurs positives comme négatives de $a$.

% BEGIN-VERIFY
% from sympy import symbols, cancel, limit
% a, h = symbols('a h', nonzero=True)
% assert cancel(((a+h)**2-a**2)/h - (2*a+h)) == 0
% assert limit(((a+h)**2-a**2)/h, h, 0) == 2*a
% assert cancel((1/(a+h)-1/a)/h + 1/(a*(a+h))) == 0
% assert limit((1/(a+h)-1/a)/h, h, 0) == -1/a**2
% # Le carré est aussi dérivable au point a=0.
% assert limit(h**2/h, h, 0) == 0
% END-VERIFY

\propriete[Extension aux exposants entiers négatifs]{
Pour tout $n \in \mathbb{Z}$ (y compris les entiers négatifs), la fonction $f(x) = x^n$ est dérivable sur $\mathbb{R} \setminus \{0\}$ et :
\[
f'(x) = nx^{n-1}.
\]
Pour $n\geqslant1$, les fonctions puissances sont aussi définies et dérivables
en $0$. Pour $n=0$, on traite la fonction constante séparément comme ci-dessus.
}

\exemple{
Pour $f(x) = x^{-2} = \dfrac{1}{x^2}$, on applique la formule avec $n = -2$ :
\[
f'(x) = -2x^{-3} = -\frac{2}{x^3}.
\]
}

\exemple{
Pour $g(x) = x^{-3} = \dfrac{1}{x^3}$, on a $n = -3$ :
\[
g'(x) = -3x^{-4} = -\frac{3}{x^4}.
\]
}

\margeAppui{La ligne $\frac{1}{x}$ du tableau des dérivées de référence est un cas particulier : $\frac{1}{x} = x^{-1}$, et la formule $(-1)x^{-2} = -\frac{1}{x^2}$ redonne bien le résultat connu.}

% BEGIN-VERIFY
% from sympy import symbols, diff, Rational
% x = symbols('x')
% # Dérivée de x^(-2) = -2*x^(-3)
% assert diff(x**(-2), x) == -2*x**(-3)
% # Dérivée de x^(-3) = -3*x^(-4)
% assert diff(x**(-3), x) == -3*x**(-4)
% # Dérivée de x^(-1) = -x^(-2) = -1/x^2
% assert diff(x**(-1), x) == -x**(-2)
% # Cas exacts supplémentaires ; cette liste ne remplace pas la preuve générale.
% for n_val in [-5, -4, -3, -2, -1, 1, 2, 3, 4, 5]:
%     assert diff(x**n_val, x) == n_val * x**(n_val - 1)
% END-VERIFY

\approfondissement{
\textbf{Puissances naturelles.} Pour $n\geqslant2$ et $a\neq0$, on peut
développer par le binôme de Newton le taux de variation
\[
\frac{(a+h)^n - a^n}{h}.
\]
Le terme constant en $h$ est $na^{n-1}$ ; tous les autres termes contiennent
une puissance strictement positive de $h$ et tendent vers $0$.
Au point $a=0$, on calcule directement, pour $h\neq0$,
\[
\frac{h^n-0}{h}=h^{n-1},
\]
qui tend vers $0$ puisque $n\geqslant2$. Le cas $n=1$ donne directement un
taux égal à $1$.

\textbf{Puissances négatives.} Pour $n=-m<0$, avec $m$ entier positif, on écrit
$x^n=\frac{1}{x^m}$ pour $x\neq0$. La règle de dérivation de l’inverse
d'une fonction, étudiée dans la capacité suivante, donne
\[
\left(\frac{1}{x^m}\right)'=-\frac{mx^{m-1}}{(x^m)^2}
=-mx^{-m-1}=nx^{n-1}.
\]
Le développement fini du binôme n'a donc été utilisé que pour un exposant naturel.
}
